\subsection{Selected General Topology}

% **Role:** abstract structure behind limits and continuity.
%
% Topological spaces, bases, product topology, quotient topology, countability axioms.


General topology abstracts the open-set behavior of metric spaces. This selective pass records spaces, bases, products, quotients, and countability axioms for later use.

\begin{definition} \label{def:topological-space} \hlt{Topological Space}\\
A \hlt{topology} on a set $X$ is a collection $\mathcal{T}\subseteq 2^X$ such that
\begin{enumerate}[label=(\roman*)]
\setlength{\itemsep}{0pt}
\item $\emptyset,X\in\mathcal{T}$;
\item arbitrary unions of members of $\mathcal{T}$ belong to $\mathcal{T}$;
\item finite intersections of members of $\mathcal{T}$ belong to $\mathcal{T}$.
\end{enumerate}
The pair $(X,\mathcal{T})$ is a \hlt{topological space}, and the members of $\mathcal{T}$ are the \hlt{open sets}. Closed sets are complements of open sets.
\end{definition}

\begin{proposition} \label{prop:metric-open-sets-form-topology} \hlt{Metric Open Sets Form a Topology}\\
If $(X,d)$ is a metric space, then the open sets from Definition \ref{def:metric-topology-core-definitions} form a topology on $X$.
\begin{proof}
The empty set and $X$ are open. Arbitrary unions and finite intersections are open by Theorem \ref{thm:metric-open-closed-set-calculus}.\\
\hlt{Structural Commentary}: Metric balls generate an open-set system with the exact axioms of a topology.\\
\hlt{Key Insight}: General topology keeps the open-set calculus and forgets the distance function.
\end{proof}
\end{proposition}

\begin{definition} \label{def:basis-for-topology} \hlt{Basis}\\
Let $X$ be a set. A collection $\mathcal{B}\subseteq 2^X$ is a \hlt{basis} for a topology if:
\begin{enumerate}[label=(\roman*)]
\setlength{\itemsep}{0pt}
\item for every $x\in X$, some $B\in\mathcal{B}$ contains $x$;
\item if $x\in B_1\cap B_2$ with $B_1,B_2\in\mathcal{B}$, then some $B_3\in\mathcal{B}$ satisfies $x\in B_3\subseteq B_1\cap B_2$.
\end{enumerate}
The topology generated by $\mathcal{B}$ consists of all unions of basis elements. A \hlt{subbasis} is a collection whose finite intersections form a basis.
\end{definition}

\begin{example} \label{ex:metric-balls-basis} \hlt{Metric Balls as a Basis}\\
In a metric space, the collection of all open balls is a basis for the metric topology. On $\R^k$, balls with rational centers and rational radii form a countable basis.
\end{example}

\begin{definition} \label{def:product-topology} \hlt{Product Topology}\\
For topological spaces $X$ and $Y$, the \hlt{product topology} on $X\times Y$ is generated by basis sets $U\times V$, where $U\subseteq X$ and $V\subseteq Y$ are open. For a family $\{X_i\}_{i\in I}$, the product topology on $\prod_iX_i$ is generated by finite-coordinate restrictions: sets of the form $\prod_iU_i$ where each $U_i$ is open and $U_i=X_i$ for all but finitely many $i$.
\end{definition}

\begin{proposition} \label{prop:product-topology-coordinate-continuity} \hlt{Product Topology and Coordinate Maps}\\
A map $f:Z\to X\times Y$ is continuous for the product topology if and only if both coordinate maps $\pi_X\circ f$ and $\pi_Y\circ f$ are continuous.
\begin{proof}
If $f$ is continuous, then coordinate maps are continuous because inverse images of open sets are inverse images of product-open cylinders. Conversely, if both coordinates are continuous, then the inverse image of each basis set $U\times V$ is $(\pi_X\circ f)^{-1}(U)\cap(\pi_Y\circ f)^{-1}(V)$, which is open.\\
\hlt{Structural Commentary}: Product topology is the weakest topology making the projections continuous.\\
\hlt{Key Insight}: Continuity into a product is checked coordinatewise.
\end{proof}
\end{proposition}

\begin{definition} \label{def:quotient-topology} \hlt{Quotient Topology}\\
Let $q:X\to Y$ be a surjective map from a topological space $X$ to a set $Y$. The \hlt{quotient topology} on $Y$ declares $V\subseteq Y$ open exactly when $q^{-1}(V)$ is open in $X$. Equivalently, it is the strongest topology on $Y$ making $q$ continuous.
\end{definition}

\begin{proposition} \label{prop:quotient-topology-universal-property} \hlt{Quotient Universal Property}\\
Let $q:X\to Y$ carry the quotient topology, and let $g:Y\to Z$ be a map into a topological space $Z$. Then $g$ is continuous if and only if $g\circ q:X\to Z$ is continuous.
\begin{proof}
If $g$ is continuous, then $g\circ q$ is continuous by composition. Conversely, if $g\circ q$ is continuous and $W\subseteq Z$ is open, then $q^{-1}(g^{-1}(W))=(g\circ q)^{-1}(W)$ is open in $X$. By the quotient topology, $g^{-1}(W)$ is open in $Y$.\\
\hlt{Structural Commentary}: Quotient topology is defined so continuity out of the quotient can be tested before identifying points.\\
\hlt{Key Insight}: Pull back through the quotient map and use the defining open-set test.
\end{proof}
\end{proposition}

\begin{definition} \label{def:topological-countability-axioms} \hlt{Countability Axioms}\\
A topological space $X$ is \hlt{first countable} if every point has a countable neighborhood basis. It is \hlt{second countable} if its topology has a countable basis. It is \hlt{separable} if it has a countable dense subset.
\end{definition}

\begin{proposition} \label{prop:second-countable-implies-separable} \hlt{Second Countable Implies Separable}\\
If $X$ is second countable, then $X$ is separable.
\begin{proof}
Let $\mathcal{B}=\{B_n:n\in\N\}$ be a countable basis. For each nonempty $B_n$, choose $x_n\in B_n$.\\
The chosen set is countable. If $U$ is a nonempty open set, some basis element $B_n$ satisfies $x\in B_n\subseteq U$, so the chosen set meets $U$. Hence it is dense.\\
\hlt{Structural Commentary}: A countable basis yields a countable dense set.\\
\hlt{Key Insight}: Density is tested by nonempty open sets.
\end{proof}
\end{proposition}
